\documentclass[
  12pt,
  openright,
  a4paper,
  twoside,
  english,
  brazil,
]{article}

\usepackage{mathtools}
\usepackage{amsfonts}
\usepackage{amssymb}
\usepackage{amsmath}
\usepackage[utf8]{inputenc}
\usepackage[portuguese]{babel}
\usepackage[T1]{fontenc}

\usepackage[top=2cm, bottom=2cm, left=2cm, right=2cm]{geometry}
\usepackage{enumerate}

\newcommand{\cqd}{\hfill $\square$}
\newcommand{\dt}[1]{
  \begin{vmatrix}#1
\end{vmatrix}}

% -----------------------------------------------------------------------------

\title{Geometria Analítica - Lista 03}
\author{Prof. Helaine - 30/04/2026}
\date{}

\begin{document}

\maketitle

\begin{enumerate}[\bfseries  1.]
  \item Determinar um vetor unitário simultaneamente ortogonal aos
    vetores $\vec u=(2,-6,3)$ e $\vec v=(4,3,1)$.

    $$\vec w=\frac{\vec u\times\vec v}{\|\vec u\times \vec v\|}$$
    $$\vec u\times\vec
    v=(\dt{-6 & 3\\3 & 1},-\dt{2 & 3\\4 & 1},\dt{2 & -6\\4 &
    3})=(-15,10,30)=5(-3,2,6)
    $$
    $$\|\vec u\times\vec v|=\sqrt{15^2+10^2+30^2}=\sqrt{1225}=35$$
    $$\vec w=\frac{1}{7}(-3,2,6)$$

  \item Dados os vetores $\vec u=(1,2,-1)$ e $\vec v=(0,-1,3)$,
    calcular a área do paralelogramo determinado pelos vetores $3\vec
    u$ e $\vec v-\vec u$.

    $$A=\|\vec u\times \vec v\|$$
    $$\vec v-\vec u=(0-1,-1-2,3-(-1))=(-1,-3,4)$$
    $$3\vec u\times\vec (\vec v-\vec u)=3(\dt{2 & -1\\-3 & 4}, -\dt{1
    & -1\\-1 & 4}, \dt{1 & 2\\-1 & -3})=3(5,-3,-1)$$
    $$A=3\sqrt{5^2+3^2+1^2}=3\sqrt{35}$$

  \item Sejam os vetores $\vec u=(3,1,-1)$ e $\vec v = (a,0,2)$.
    Calcular o valor de $a$ para que a área do paralelogramo
    determinado por $\vec u$ e $\vec v$ seja igual a $2\sqrt 6$.

    $$
    \vec u\times\vec v=(\dt{1 & -1\\0 & 2}, -\dt{3 & -1\\a & 2}, \dt{3
    & 1\\a & 0})=(2,-(6+a),-a)
    $$
    $$\|\vec u\times\vec v\|=\sqrt{2^2+(6+a)^2+(-a)^2}=2\sqrt{6}$$
    $$4+(6+a-a)^2-2(6+a)(-a)=4\cdot 6$$
    $$4+36+12a+2a^2-24=0$$
    $$2a^2+12a+16=0$$
    $$a^2+6a+8=0$$
    $$a=
    \begin{cases}-2\\-4
    \end{cases}$$

\end{enumerate}

\end{document}
